\documentclass[../cours_spe_trigo.tex]{subfiles}
\begin{document}
\setcounter{chapter}{0}
\chapter{D'où vient \texorpdfstring{$\pi$}{pi} ? Du cercle à sa circonférence}
\section{Les objets géométriques de départ}
\begin{definition}
Un \textbf{cercle} de centre $O$ et de rayon $R>0$ est l'ensemble des points du plan situés à la distance $R$ de $O$. Le \textbf{disque} est l'ensemble des points situés à une distance au plus $R$. Le \textbf{diamètre} vaut $D=2R$. La \textbf{circonférence} $C$ désigne la longueur du contour, pas l'aire du disque.
\end{definition}
\begin{center}
\begin{tikzpicture}[scale=1]
\fill[bleuPale] (0,0) circle (1.65);
\draw[very thick,bleu] (0,0) circle (1.65);
\draw[thick,orange] (-1.65,0)--(1.65,0);
\draw[thick,bleuNuit] (0,0)--(48:1.65);
\fill (0,0) circle (.035);\node[above left] at (0,0) {$O$};
\node[above left,bleuNuit] at (24:1.0) {$R$};
\node[below,orange] at (0,-.25) {$D=2R$};
\node[bleu,right] at (1.7,.8) {cercle};
\node[bleuNuit] at (.4,-1.0) {disque};
\end{tikzpicture}
\end{center}
\textbf{Prérequis sur les longueurs.} Deux figures obtenues l'une de l'autre par une homothétie de rapport $\lambda>0$ ont toutes leurs longueurs multipliées par $\lambda$ et toutes leurs aires multipliées par $\lambda^2$. Ce résultat de similitude, connu pour les segments, reste vrai pour les courbes lorsque leur longueur est définie par approximation polygonale.

\section{Pourquoi le quotient circonférence/diamètre est-il constant ?}
Ce qui mérite une explication n'est \emph{pas} l'égalité $C=\pi D$, qui découlera d'une définition, mais le fait que \textbf{tous les cercles ont le même rapport $C/D$}.
\begin{demonstration}
\textbf{Démonstration par homothétie.} Soient deux cercles de rayons $R_1$ et $R_2$. L'homothétie de centre $O$ et de rapport $\lambda=R_2/R_1$ transforme le premier cercle en un cercle de rayon $R_2$. Elle multiplie la longueur de tout segment par $\lambda$. En approchant les deux circonférences par des polygones correspondants, elle multiplie également les périmètres de ces polygones par $\lambda$. En passant à la limite qui définit la longueur de la courbe, on obtient
\[
C_2=\lambda C_1=\frac{R_2}{R_1}C_1.
\]
Par ailleurs $D_2=\lambda D_1$. En divisant, $C_2/D_2=C_1/D_1$. Le quotient ne dépend donc pas du cercle. $\square$
\end{demonstration}
\begin{definition}
On \textbf{appelle $\pi$} le rapport commun entre la longueur d'un cercle et son diamètre :
\[\boxed{\displaystyle\pi=\frac{C}{D}}\]
Ainsi $C=\pi D$, et comme $D=2R$ :
\[\boxed{\displaystyle C=2\pi R.}\]
Il ne faut pas dire « le périmètre d'un cercle vaut $2\pi$ » sans préciser \emph{le rayon $R=1$}.
\end{definition}
\textbf{À propos de la valeur de $\pi$.} Le nombre $\pi$ est une constante géométrique, approximativement $3{,}1415926535\ldots$. Son écriture décimale ne s'arrête pas et n'est pas périodique : $\pi$ est irrationnel (preuve hors programme du lycée). La formule $C=2\pi R$ est \emph{exacte} ; $C\approx 6{,}2832R$ est une approximation.

\section{Comment calculer $\pi$ sans connaître d'avance sa valeur ?}
Une méthode ancienne consiste à \textbf{encadrer} le cercle entre deux polygones réguliers : l'un inscrit dans le cercle, l'autre circonscrit au cercle.
\begin{center}
\begin{tikzpicture}[scale=1.0]
\begin{scope}[xshift=-3cm]
\draw[bleu,thick] (0,0) circle (1.3);
\draw[orange,very thick] (0:1.3)--(60:1.3)--(120:1.3)--(180:1.3)--(240:1.3)--(300:1.3)--cycle;
\node[below] at (0,-1.65) {hexagone inscrit};
\end{scope}
\begin{scope}[xshift=3cm]
\draw[bleu,thick] (0,0) circle (1.3);
\foreach \k in {0,...,5} {\draw[vert!75!black,very thick] ({30+60*\k}:1.5011)--({90+60*\k}:1.5011);}
\node[below] at (0,-1.65) {hexagone circonscrit};
\end{scope}
\end{tikzpicture}
\end{center}
\textbf{Exemple : l'hexagone, avec $R=1$.} Un hexagone régulier inscrit est formé de six triangles équilatéraux de côté $1$, donc son périmètre vaut $P_6=6$. Dans l'hexagone circonscrit, chaque demi-côté mesure $\tan(30^\circ)=1/\sqrt3$ (ou, sans trigonométrie, par Pythagore), donc $Q_6=12/\sqrt3=4\sqrt3$. On en déduit l'encadrement
\[
6<C(1)<4\sqrt3\qquad\Longrightarrow\qquad 3<\pi<2\sqrt3.
\]
\textbf{En augmentant le nombre de côtés, l'encadrement devient précis :}
\begin{center}
\begin{tabular}{rcc}
\toprule
Nombre de côtés $n$ & $P_n/2$ (borne inférieure) & $Q_n/2$ (borne supérieure) \\
\midrule
$6$ & $3{,}000000$ & $3{,}464102$\\
$12$ & $3{,}105829$ & $3{,}215390$\\
$24$ & $3{,}132629$ & $3{,}159660$\\
$48$ & $3{,}139350$ & $3{,}146086$\\
$96$ & $3{,}141032$ & $3{,}142715$\\
\bottomrule
\end{tabular}
\end{center}
\begin{demonstration}
\textbf{Pourquoi les deux suites convergent-elles vers une même longueur ?} Si $c_n$ est le côté d'un polygone régulier inscrit dans un cercle de rayon $1$ et $a_n$ son apothème (distance du centre à un côté), Pythagore donne $a_n=\sqrt{1-c_n^2/4}$, $P_n=nc_n$ et $Q_n=P_n/a_n$. En doublant le nombre de côtés, $c_n$ diminue et $a_n\to1$, donc $Q_n/P_n=1/a_n\to1$. Les périmètres inscrits croissent, les circonscrits décroissent ; ils encadrent une même longueur-limite $C(1)$. Cette \emph{rectification} du cercle justifie géométriquement l'existence d'une circonférence. Le calcul des premières valeurs ne suppose aucune valeur préalable de $\pi$ : on obtient les nouveaux côtés avec Pythagore.
\end{demonstration}
\begin{attention}
On ne \textbf{démontre pas} la formule $C=2\pi R$ en remplaçant simplement $\pi$ par $3{,}14$ : la démonstration consiste à établir la constance de $C/D$, puis à \emph{définir} cette constante comme $\pi$. L'approximation numérique intervient ensuite.
\end{attention}

\section{Pourquoi l'aire du disque est-elle $\pi R^2$ ?}
\begin{propriete}
L'aire d'un disque de rayon $R$ est $\boxed{\mathcal A=\pi R^2}$.
\end{propriete}
\begin{demonstration}
\textbf{Démonstration par polygones réguliers.} Un polygone régulier inscrit à $n$ côtés est une réunion de $n$ triangles isocèles de base $c_n$ et de hauteur $a_n$. Son aire est donc $A_n=n\cdot\frac{c_na_n}{2}=\frac12P_na_n$. Quand $n$ augmente, $P_n\to C=2\pi R$ et $a_n\to R$. Les aires de ces polygones approchent l'aire du disque ; en passant à la limite,
\[
\mathcal A=\lim_{n\to\infty}\frac12P_na_n=\frac12(2\pi R)R=\boxed{\pi R^2}.
\]
Cette preuve utilise le fait qu'un disque peut être approché par l'intérieur avec une précision arbitraire par des polygones réguliers. $\square$
\end{demonstration}
\textbf{Exemple.} Si $R=3\,\mathrm{cm}$, alors $C=6\pi\,\mathrm{cm}$ (une longueur) et $\mathcal A=9\pi\,\mathrm{cm}^2$ (une aire) : les unités ne sont pas interchangeables.
\end{document}
